\documentclass[10pt,a4paper]{amsart}
\usepackage[margin=19mm]{geometry}
\usepackage{amsmath,amssymb,mathtools}
\usepackage[scr=rsfs,cal=euler]{mathalfa}
\usepackage{xfrac}
\usepackage[unicode,hidelinks,bookmarks=false]{hyperref}
\newcommand{\ie}{i.e.\ }
\newcommand{\cf}{cf.\ }
\newcommand{\resp}{respectively\ }
\newcommand{\Subsection}{\subsection}
\providecommand{\nobreakdash}{\nobreak\hbox{-}}
\setlength{\emergencystretch}{2em}
\multlinegap0pt
\usepackage{amssymb}
\usepackage{mathtools}
\usepackage[shortlabels]{enumitem}
\usepackage{xcolor}
\newtheorem{lemma}{Lemma}
\newcommand{\TMF}{\mathrm{TMF}}
\title{The hyperbolic class and vanishing laws\\in a TMF-valued four-manifold invariant}
\author{Yuqi Li, Hao-Yu Sun}
\date{Source: arXiv:2609.09819v1 (CC BY 4.0)}
\begin{document}
\maketitle

\noindent\emph{Source and licence.} The hyperbolic class and vanishing laws\\in a TMF-valued four-manifold invariant. Version arXiv:2609.09819v1 (CC BY 4.0), distributed by arXiv under the Creative Commons Attribution 4.0 International licence, http://creativecommons.org/licenses/by/4.0/, as recorded in the arXiv licence metadata for this exact version and shown on the abstract page https://arxiv.org/abs/2609.09819v1. Full licence notice: \texttt{licenses/arxiv-2609.09819-CC-BY-4.0.txt}.\par\smallskip

\noindent\textit{Editorial scope.} Counted unit: the article's lemma lem:H-shift (source0.tex lines 545--550). The article's complete original proof of it is delivered with it (source0.tex lines 552--579, one \texttt{proof} environment of 28 lines). No mathematics is written, repaired or re-derived: the statement and the proof are the article's own lines, in the article's own notation, and no retained line is substituted or re-flowed.  No further source-local context is needed: every cross-reference of every retained line resolves inside the two delivered spans.  Nothing was omitted from any retained span, so the package carries no omission record.  No retained line contains a citation, so the wrapper carries no bibliography and no bibliography entry.  No retained line contains a figure environment, an image-inclusion command or drawing code, so no image is delivered and the package declares no asset dependency.  Editorial formatting only, in addition: an AMS \texttt{amsart} preamble that restates the article's own declarations and macros used by the retained lines, namely the environments \texttt{lemma} on separate counters, together with the standard packages those lines need.  The retained environments print in this document as Lemma 1 in order of appearance, whereas the article prints the same items with the numbering of its own full document.  The complete native source of the article as distributed in the arXiv e-print for this exact version is bundled unchanged as \texttt{sources/209853/source0.tex}.
\begin{lemma}\label{lem:H-shift}
There is an equivalence
\[
 L^H\simeq\TMF[1].
\]
\end{lemma}

\begin{proof}
The canonical rank-one shift gives
\[
 L^H[-2]\simeq L^{H\oplus\langle1\rangle}.
\]
Set
\[
 P=\begin{pmatrix}
 1&0&1\\
 0&1&-1\\
 1&-1&1
 \end{pmatrix}.
\]
The matrix satisfies $\det P=-1$ and
\[
 P(H\oplus\langle1\rangle)P^{\mathsf T}
 =\operatorname{diag}(1,1,-1).
\]
Applying the rank-one shifts to the resulting diagonal form gives
\[
 L^{H\oplus\langle1\rangle}
 \simeq
 L^{\langle1\rangle}\otimes L^{\langle1\rangle}\otimes L^{\langle-1\rangle}
 \simeq\TMF[-1].
\]
Hence $L^H[-2]\simeq\TMF[-1]$, which is equivalent to the
claimed shift.
\end{proof}

\end{document}
